%% =========================================================
%%  AUTH & SESSION SECURITY ANALYZER
%%  Rapport Technique Complet — Architecture, Outils & Pipeline
%%  Compilez avec: pdflatex rapport.tex (2 passes)
%% =========================================================
\documentclass[12pt,a4paper,french]{article}

%% ── Encodage & Langue ──────────────────────────────────────
\usepackage[utf8]{inputenc}
\usepackage[T1]{fontenc}
\usepackage[french]{babel}
\usepackage{csquotes}

%% ── Mise en page ────────────────────────────────────────────
\usepackage[top=2.5cm, bottom=2.5cm, left=3cm, right=3cm]{geometry}
\usepackage{setspace}
\setstretch{1.15}

%% ── Typographie ─────────────────────────────────────────────
\usepackage{lmodern}
\usepackage{microtype}

%% ── Couleurs ────────────────────────────────────────────────
\usepackage{xcolor}
\definecolor{primary}{HTML}{3B82F6}       % bleu
\definecolor{purple}{HTML}{8B5CF6}        % violet
\definecolor{green}{HTML}{10B981}         % vert
\definecolor{red}{HTML}{EF4444}           % rouge
\definecolor{orange}{HTML}{F59E0B}        % orange
\definecolor{darkbg}{HTML}{0F172A}        % fond sombre
\definecolor{codebg}{HTML}{1E293B}        % fond code
\definecolor{mutedtxt}{HTML}{64748B}      % texte secondaire
\definecolor{mavscolor}{HTML}{6366F1}     % indigo MASVS

%% ── Graphiques & TikZ ───────────────────────────────────────
\usepackage{tikz}
\usetikzlibrary{shapes.geometric, arrows.meta, positioning, calc,
                decorations.pathreplacing, backgrounds, fit, shadows}
\usepackage{pgfplots}
\pgfplotsset{compat=1.18}

%% ── Listes & Tableaux ───────────────────────────────────────
\usepackage{enumitem}
\usepackage{booktabs}
\usepackage{tabularx}
\usepackage{longtable}
\usepackage{multirow}
\usepackage{array}
\usepackage{colortbl}

%% ── Code source ─────────────────────────────────────────────
\usepackage{listings}
\lstset{
  basicstyle=\ttfamily\small,
  backgroundcolor=\color{codebg},
  rulecolor=\color{primary},
  frame=single,
  framesep=4pt,
  framerule=0.4pt,
  xleftmargin=10pt,
  xrightmargin=10pt,
  breaklines=true,
  breakatwhitespace=true,
  showstringspaces=false,
  keywordstyle=\color{primary}\bfseries,
  commentstyle=\color{mutedtxt}\itshape,
  stringstyle=\color{green},
  numberstyle=\tiny\color{mutedtxt},
  numbers=left,
  numbersep=8pt,
  columns=flexible,
  keepspaces=true,
  literate={é}{{\'{e}}}1 {è}{{\`{e}}}1 {à}{{\`{a}}}1 {ê}{{\^{e}}}1
           {î}{{\^{i}}}1 {ô}{{\^{o}}}1 {û}{{\^{u}}}1 {ù}{{\`{u}}}1
           {â}{{\^{a}}}1 {ç}{{\c{c}}}1 {É}{{\'{E}}}1 {À}{{\`{A}}}1
}

%% ── Boîtes colorées ─────────────────────────────────────────
\usepackage[most]{tcolorbox}
\tcbuselibrary{listings, breakable, skins}

\tcbset{
  mavsbox/.style={
    enhanced, breakable,
    colback=mavscolor!8!white,
    colframe=mavscolor,
    fonttitle=\bfseries\small,
    coltitle=white,
    attach boxed title to top left={yshift=-2mm, xshift=4mm},
    boxed title style={colback=mavscolor, rounded corners},
    arc=4pt, boxrule=0.6pt,
    before skip=8pt, after skip=8pt
  },
  toolbox/.style={
    enhanced,
    colback=primary!6!white,
    colframe=primary!60!white,
    fonttitle=\bfseries\small,
    arc=4pt, boxrule=0.6pt,
    before skip=6pt, after skip=6pt
  },
  warnbox/.style={
    enhanced,
    colback=orange!8!white,
    colframe=orange!70!white,
    fonttitle=\bfseries\small,
    arc=4pt, boxrule=0.6pt,
    before skip=6pt, after skip=6pt
  },
  successbox/.style={
    enhanced,
    colback=green!6!white,
    colframe=green!60!white,
    fonttitle=\bfseries\small,
    arc=4pt, boxrule=0.6pt,
    before skip=6pt, after skip=6pt
  },
  dangerbox/.style={
    enhanced,
    colback=red!6!white,
    colframe=red!60!white,
    fonttitle=\bfseries\small\color{red},
    arc=4pt, boxrule=0.6pt,
    before skip=6pt, after skip=6pt
  }
}

%% ── Liens hypertexte ────────────────────────────────────────
\usepackage{hyperref}
\hypersetup{
  colorlinks=true,
  linkcolor=primary,
  urlcolor=purple,
  citecolor=green,
  pdftitle={Auth \& Session Security Analyzer — Rapport Technique},
  pdfauthor={Projet Antigravity},
  pdfsubject={OWASP MASVS · DevSecOps · Mobile AppSec}
}

%% ── En-têtes & Pieds de page ────────────────────────────────
\usepackage{fancyhdr}
\pagestyle{fancy}
\fancyhf{}
\fancyhead[L]{\small\textcolor{mutedtxt}{Auth \& Session Security Analyzer}}
\fancyhead[R]{\small\textcolor{mutedtxt}{Rapport Technique}}
\fancyfoot[C]{\small\textcolor{mutedtxt}{\thepage}}
\renewcommand{\headrulewidth}{0.4pt}
\renewcommand{\headrule}{\hbox to\headwidth{\color{primary}\leaders\hrule height\headrulewidth\hfill}}

%% ── Sections personnalisées ─────────────────────────────────
\usepackage{titlesec}
\titleformat{\section}
  {\color{primary}\large\bfseries}{\thesection}{0.8em}{}
  [\color{primary}\titlerule]
\titleformat{\subsection}
  {\color{purple}\normalsize\bfseries}{\thesubsection}{0.7em}{}
\titleformat{\subsubsection}
  {\color{mutedtxt}\normalsize\bfseries}{\thesubsubsection}{0.6em}{}

%% ── Tables des matières ─────────────────────────────────────
\usepackage{tocloft}
\renewcommand{\cftsecfont}{\color{primary}\bfseries}
\renewcommand{\cftsubsecfont}{\color{purple}}
\renewcommand{\cftsecpagefont}{\color{primary}}

%% ── Commandes utilitaires ───────────────────────────────────
\newcommand{\badge}[2]{%
  \colorbox{#1!15}{\textcolor{#1}{\small\textbf{#2}}}%
}
\newcommand{\masvs}[1]{\badge{mavscolor}{\texttt{#1}}}
\newcommand{\mastg}[1]{\badge{primary}{\texttt{#1}}}
\newcommand{\tool}[1]{\textcolor{purple}{\textbf{#1}}}
\newcommand{\file}[1]{\textcolor{primary}{\texttt{#1}}}
\newcommand{\critical}[1]{\textcolor{red}{\textbf{#1}}}
\newcommand{\ok}[1]{\textcolor{green}{\textbf{#1}}}

%% ── Métadonnées ─────────────────────────────────────────────
\title{%
  \vspace{-1cm}
  {\Huge\textbf{\textcolor{primary}{Auth \& Session}}}\\[4pt]
  {\Huge\textbf{\textcolor{purple}{Security Analyzer}}}\\[8pt]
  {\large\textcolor{mutedtxt}{Rapport Technique Complet}}\\[4pt]
  {\normalsize\textcolor{mutedtxt}{Architecture · Outils · Pipeline d'Analyse · Standards OWASP MASVS v2}}
}
\author{%
  \textbf{Projet Antigravity}\\
  \small\textcolor{mutedtxt}{DevSecOps · Mobile AppSec · OWASP MASVS/MASTG}
}
\date{\textcolor{mutedtxt}{Avril 2026}}

%% =========================================================
\begin{document}
%% =========================================================

\maketitle
\thispagestyle{empty}

\vspace{0.5cm}

%% Bannière descriptive
\begin{tcolorbox}[
  enhanced,
  colback=darkbg,
  colframe=primary,
  fontupper=\color{white}\normalsize,
  arc=8pt, boxrule=1pt,
  left=12pt, right=12pt, top=10pt, bottom=10pt
]
\textbf{Auth \& Session Security Analyzer} est un outil d'audit automatisé de la sécurité
d'authentification pour applications mobiles Android.
Il combine l'analyse \textbf{SAST} (statique), l'analyse \textbf{DAST} (dynamique),
une intelligence artificielle \textbf{RAG + Gemini} et une couverture complète des standards
\textbf{OWASP MASVS v2} / \textbf{MASTG} pour produire des Checklists de sécurité,
des critères d'acceptation \textbf{Agile} et des scénarios \textbf{Gherkin} prêts pour Jira.
\end{tcolorbox}

\vspace{0.4cm}
\tableofcontents
\newpage

%% ==========================================================
\section{Présentation Générale du Projet}
%% ==========================================================

\subsection{Contexte et Objectifs}

Les applications mobiles exposent souvent des vulnérabilités critiques liées à
l'authentification et à la gestion des sessions : tokens mal stockés, sessions
non révoquées après déconnexion, algorithmes faibles ou absents dans les JWT,
absence de certificate pinning, etc.

L'objectif de l'\textbf{Auth \& Session Security Analyzer} est de résoudre ce problème
en agissant comme un \textit{auditeur DevSecOps automatisé} qui :

\begin{itemize}[leftmargin=1.5cm, itemsep=4pt]
  \item[\textcolor{primary}{$\blacktriangleright$}]
    Analyse les APKs Android (analyse statique SAST) et les logs de trafic réseau (DAST)
  \item[\textcolor{primary}{$\blacktriangleright$}]
    Génère des checklists de sécurité alignées \textbf{MASVS v2} grâce à l'IA
  \item[\textcolor{primary}{$\blacktriangleright$}]
    Produit des \textit{Security Acceptance Criteria} directement intégrables dans des User Stories Jira
  \item[\textcolor{primary}{$\blacktriangleright$}]
    Implémente le principe \textbf{Shift-Left} : la sécurité est intégrée dès la conception
  \item[\textcolor{primary}{$\blacktriangleright$}]
    Exporte des rapports complets en Markdown et JSON pour les équipes sécurité
\end{itemize}

\subsection{Stack Technologique}

\begin{center}
\renewcommand{\arraystretch}{1.4}
\begin{tabular}{lll}
\toprule
\textbf{Couche} & \textbf{Technologies} & \textbf{Version} \\
\midrule
\rowcolor{primary!6}
Backend API     & Python, FastAPI, Uvicorn            & Python 3.11+ / FastAPI 0.110+ \\
Base de données & SQLite, SQLAlchemy ORM              & SQLAlchemy 2.0+ \\
\rowcolor{primary!6}
IA / RAG        & Google Gemini, LangChain, ChromaDB  & Gemini 2.5 Flash \\
Embeddings      & SentenceTransformer (MiniLM)        & all-MiniLM-L6-v2 \\
\rowcolor{primary!6}
SAST            & Semgrep, Apktool, JADX              & Semgrep 1.66+ \\
DAST            & Parsers custom Python               & PyJWT 2.8+ \\
\rowcolor{primary!6}
Frontend        & React, TypeScript, TailwindCSS, Vite & React 18 / Vite 5 \\
\bottomrule
\end{tabular}
\end{center}

%% ==========================================================
\section{Architecture du Projet}
%% ==========================================================

\subsection{Vue d'Ensemble de la Structure}

\begin{lstlisting}[language=bash, caption={Structure des répertoires du projet}]
Auth-and-Session-Security-Analyzer/
├── backend/
│   ├── main.py                  # Point d'entrée FastAPI — registration des routes
│   ├── core/
│   │   ├── config.py            # Paramètres (Pydantic Settings, .env)
│   │   ├── database.py          # Modèles SQLAlchemy + SessionLocal
│   │   └── models.py            # Modèles Pydantic (Finding, AuditReport, etc.)
│   ├── api/routes/
│   │   ├── upload.py            # POST /api/upload/apk — réception de fichiers
│   │   ├── analyze.py           # POST /api/analyze/start — lancement du pipeline
│   │   ├── jobs.py              # GET  /api/jobs/ — suivi des jobs en cours
│   │   ├── report.py            # GET  /api/report/ — consultation des rapports
│   │   ├── export.py            # GET  /api/export/{id}/markdown|json
│   │   └── auth_audit.py        # POST /api/auth-audit/generate-checklist
│   ├── ai/
│   │   ├── llm_client.py        # Client Gemini via LangChain
│   │   └── rag_engine.py        # Moteur RAG ChromaDB + SentenceTransformer
│   ├── analyzers/
│   │   ├── runner.py            # Chef d'orchestre du pipeline d'analyse
│   │   ├── semgrep_runner.py    # Wrapper Semgrep (SAST)
│   │   └── flow_analyzer.py     # Analyse du trafic HTTP (DAST)
│   ├── parsers/
│   │   ├── apk_parser.py        # Décompilation APK (Apktool + JADX)
│   │   ├── manifest_parser.py   # Lecture AndroidManifest.xml
│   │   ├── traffic_parser.py    # Analyse logs HTTP/Burp
│   │   └── jwt_parser.py        # Décodage et audit des tokens JWT
│   └── rules/
│       ├── semgrep/auth.yaml    # 7 règles SAST custom MASVS-alignées
│       └── masvs/masvs_auth_full.json  # 17 règles MASVS base RAG
├── frontend/src/
│   ├── App.tsx                  # Routage par onglets
│   ├── components/
│   │   ├── Scanner.tsx          # Interface upload + suivi analyse
│   │   ├── AuthAudit.tsx        # Formulaire d'audit IA + résultats
│   │   ├── MASVSMapping.tsx     # Cartographie des contrôles MASVS
│   │   └── ReportHistory.tsx    # Historique jobs + rapports
│   └── index.css                # Système de design dark (variables CSS)
├── .env                         # Configuration locale (secrets)
└── docker-compose.yml           # Déploiement conteneurisé
\end{lstlisting}

\subsection{Diagramme d'Architecture Logicielle}

\begin{figure}[ht]
\centering
\begin{tikzpicture}[
  node distance=1.2cm and 1.8cm,
  every node/.style={font=\small},
  layer/.style={draw=primary, fill=primary!8, rounded corners=6pt,
                text width=3.5cm, align=center, minimum height=0.9cm,
                drop shadow={opacity=0.2}},
  component/.style={draw=#1!60, fill=#1!10, rounded corners=4pt,
                    text width=2.8cm, align=center, minimum height=0.8cm,
                    font=\scriptsize\bfseries},
  arrow/.style={-Stealth, thick, color=primary!70},
  label/.style={font=\scriptsize\itshape, color=mutedtxt}
]

%% Ligne 1 — Frontend
\node[layer] (fe) {\textbf{\textcolor{primary}{Frontend}}\\\scriptsize React + TypeScript};

%% Ligne 2 — API FastAPI
\node[layer, below=1.4cm of fe] (api) {\textbf{\textcolor{purple}{API FastAPI}}\\\scriptsize REST / JSON};

%% Ligne 3 — Modules backend (côte à côte)
\node[component=green, below left=1.3cm and 2.8cm of api] (upload) {Upload\\\texttt{/api/upload}};
\node[component=orange, below left=1.3cm and 0.5cm of api] (analyze) {Analyze\\\texttt{/api/analyze}};
\node[component=purple, below right=1.3cm and 0.5cm of api] (audit) {Auth Audit\\\texttt{/api/auth-audit}};
\node[component=primary, below right=1.3cm and 2.8cm of api] (export) {Export\\\texttt{/api/export}};

%% Ligne 4 — Pipeline d'analyse
\node[component=orange, below=1.4cm of analyze] (runner) {\textbf{Runner}\\\scriptsize Orchestrateur};

%% Ligne 4 — Parsers
\node[component=primary, below left=1.3cm and 1cm of runner] (apk) {ApkParser\\\scriptsize Apktool+JADX};
\node[component=primary, below right=1.3cm and 1cm of runner] (traffic) {TrafficParser\\\scriptsize Burp/HAR};

%% Ligne 5 — Analyzers
\node[component=red, below=1.2cm of apk] (semgrep) {Semgrep\\\scriptsize SAST Rules};
\node[component=orange, below=1.2cm of traffic] (flow) {FlowAnalyzer\\\scriptsize DAST HTTP};
\node[component=purple, below=1.2cm of runner] (jwt) {JWTParser\\\scriptsize Token Audit};

%% Ligne 6 — IA
\node[component=purple, below left=1.3cm and 0.2cm of jwt] (rag) {RAG Engine\\\scriptsize ChromaDB};
\node[component=primary, below right=1.3cm and 0.2cm of jwt] (llm) {LLM Client\\\scriptsize Gemini API};

%% Ligne 7 — Base
\node[component=green, below=1.2cm of rag] (db) {SQLite DB\\\scriptsize SQLAlchemy};
\node[component=orange, below=1.2cm of llm] (mavsdb) {MASVS KB\\\scriptsize 17 règles JSON};

%% Flèches
\draw[arrow] (fe) -- (api) node[label, midway, right=2pt] {HTTP/REST};
\draw[arrow] (api) -- (upload);
\draw[arrow] (api) -- (analyze);
\draw[arrow] (api) -- (audit);
\draw[arrow] (api) -- (export);
\draw[arrow] (analyze) -- (runner);
\draw[arrow] (runner) -- (apk);
\draw[arrow] (runner) -- (traffic);
\draw[arrow] (apk) -- (semgrep);
\draw[arrow] (traffic) -- (flow);
\draw[arrow] (runner) -- (jwt);
\draw[arrow] (runner) -- (rag);
\draw[arrow] (runner) -- (llm);
\draw[arrow] (rag) -- (db);
\draw[arrow] (rag) -- (mavsdb);
\draw[arrow] (llm.south) to[bend right=10] (db.east);

\end{tikzpicture}
\caption{Architecture logicielle de l'Auth \& Session Security Analyzer}
\end{figure}

%% ==========================================================
\section{Description Détaillée de Chaque Outil}
%% ==========================================================

\subsection{Backend Core — \file{backend/core/}}

\subsubsection{\file{config.py} — Gestion de la Configuration}

\begin{tcolorbox}[toolbox, title={Rôle : Centralisation de la Configuration}]
\textbf{Technologie :} Pydantic Settings + python-dotenv\\
\textbf{Responsabilité :} Lecture et validation de toutes les variables d'environnement.
La classe \texttt{Settings} valide automatiquement les types, génère des chemins absolus
cross-platform (Windows / Linux), et crée les répertoires nécessaires au démarrage.
\end{tcolorbox}

Variables principales gérées :

\begin{center}
\renewcommand{\arraystretch}{1.3}
\begin{tabular}{llp{6cm}}
\toprule
\textbf{Variable} & \textbf{Défaut} & \textbf{Rôle} \\
\midrule
\rowcolor{primary!5}
\texttt{GEMINI\_API\_KEY} & \textit{None} & Clé Google Gemini AI (RAG + Checklist) \\
\texttt{DATABASE\_URL} & SQLite local & Connexion base de données \\
\rowcolor{primary!5}
\texttt{UPLOAD\_DIR} & \texttt{data/uploads/} & Stockage temporaire des fichiers uploadés \\
\texttt{CHROMA\_DIR} & \texttt{data/chroma\_db/} & Base vectorielle ChromaDB persistante \\
\rowcolor{primary!5}
\texttt{MOBSF\_URL} & \textit{None} & URL MobSF (\textbf{port 8008}, pas 8000) \\
\texttt{MOBSF\_API\_KEY} & \textit{None} & Clé API MobSF (depuis le dashboard) \\
\bottomrule
\end{tabular}
\end{center}

\begin{tcolorbox}[warnbox, title={Attention — Conflit de ports MobSF}]
MobSF ne doit \textbf{jamais} être configuré sur le port 8000 (réservé à FastAPI).
Utiliser le port 8008 : \texttt{docker run -p 8008:8000 opensecurity/mobile-security-framework-mobsf}\\
API Key MobSF : accessible sur \texttt{http://localhost:8008} $\rightarrow$ Menu $\rightarrow$ REST API.
\end{tcolorbox}

\subsubsection{\file{database.py} — Couche de Persistance}

\begin{tcolorbox}[toolbox, title={Rôle : ORM SQLite + Modèles de Base de Données}]
\textbf{Technologie :} SQLAlchemy 2.0, SQLite\\
\textbf{Responsabilité :} Définit deux tables principales et expose une fabrique
de sessions (\texttt{SessionLocal}).
\end{tcolorbox}

\begin{center}
\renewcommand{\arraystretch}{1.3}
\begin{tabular}{lp{9cm}}
\toprule
\textbf{Table} & \textbf{Description} \\
\midrule
\rowcolor{primary!5}
\texttt{reports} & Stocke le rapport JSON complet de chaque analyse (APK name, hash, JSON sérialisé) \\
\texttt{analysis\_jobs} & Suit l'état d'un job : \texttt{pending} → \texttt{running} → \texttt{done} / \texttt{failed} \\
\bottomrule
\end{tabular}
\end{center}

\textbf{Pattern SessionLocal :} Pour éviter les sessions partagées entre threads (bug critique
dans les BackgroundTasks), chaque opération DB crée sa propre session via \texttt{SessionLocal()}
et la ferme dans un bloc \texttt{try/finally}.

\subsubsection{\file{models.py} — Modèles de Données Métier}

\begin{tcolorbox}[toolbox, title={Rôle : Contrats de Données Pydantic}]
\textbf{Technologie :} Pydantic v2 (avec \texttt{@computed\_field})\\
\textbf{Responsabilité :} Définit la structure de toutes les données métier.
\end{tcolorbox}

\begin{description}[leftmargin=2cm, style=nextline]
  \item[\texttt{Finding}]
    Représente une vulnérabilité détectée. Contient : titre, description, fichier, ligne,
    preuve (\textit{evidence}), identifiant MASVS, test MASTG, et les scores
    \texttt{impact} × \texttt{exploitability} × \texttt{exposure} (1-3 chacun, max 27).
    La \textbf{sévérité} est calculée automatiquement : critique ($\geq 18$), majeur ($\geq 9$), mineur.

  \item[\texttt{AuditReport}]
    Regroupe tous les \texttt{Finding}s d'une analyse. Calcule le \textbf{score global}
    sur 100 en déduisant des pénalités par sévérité (15 pts/critique, 5 pts/majeur, 1 pt/mineur).

  \item[\texttt{ChecklistItem} / \texttt{SACItem}]
    Représentent respectivement un élément de checklist MASVS et un critère d'acceptation Agile
    (Security Acceptance Criteria).
\end{description}

%% ----------------------------------------------------------
\subsection{Parseurs — \file{backend/parsers/}}
%% ----------------------------------------------------------

\subsubsection{\file{apk\_parser.py} — Décompilation des APKs Android}

\begin{tcolorbox}[toolbox, title={Rôle : Extraction du contenu d'un fichier APK}]
\textbf{Outils externes :} Apktool, JADX\\
\textbf{Responsabilité :} Décompile l'APK Android en deux étapes successives.
\end{tcolorbox}

\begin{enumerate}[leftmargin=1.8cm, itemsep=4pt]
  \item \textbf{Apktool} (\texttt{run\_apktool()}) : décompresse l'APK, décode les ressources
    XML binaires (notamment \texttt{AndroidManifest.xml}) et les fichiers \texttt{smali}
    (bytecode Dalvik). Sortie dans \texttt{extract\_dir/\{hash\}/apktool/}.

  \item \textbf{JADX} (\texttt{run\_jadx()}) : décompile le bytecode \texttt{.dex}
    en code Java lisible. Sortie dans \texttt{extract\_dir/\{hash\}/jadx/}.
    Ce répertoire est ensuite passé à \tool{Semgrep} pour l'analyse SAST.
\end{enumerate}

\textbf{Cache intégré :} Si le répertoire de sortie existe déjà, la décompilation est
ignorée (optimisation pour les analyses répétées du même APK).

\subsubsection{\file{manifest\_parser.py} — Audit du AndroidManifest}

\begin{tcolorbox}[toolbox, title={Rôle : Analyse de sécurité du Manifest Android}]
\textbf{Technologie :} \texttt{xml.etree.ElementTree} (bibliothèque standard Python)\\
\textbf{Responsabilité :} Analyse le fichier XML décodé pour détecter des misconfigurations
de sécurité Android.
\end{tcolorbox}

Détections réalisées :

\begin{itemize}[leftmargin=1.5cm]
  \item \textbf{Activités exportées} (\texttt{android:exported=true}) sans permission requise
    $\rightarrow$ \mavscs{MASVS-AUTH-2}
  \item \textbf{Permissions dangereuses} (CAMERA, READ\_CONTACTS, RECORD\_AUDIO)
    $\rightarrow$ \mavscs{MASVS-AUTH-1}
\end{itemize}

\subsubsection{\file{traffic\_parser.py} — Analyse des Logs de Trafic HTTP}

\begin{tcolorbox}[toolbox, title={Rôle : Extraction d'artefacts de sécurité depuis les captures réseau}]
\textbf{Formats supportés :} Logs Burp Suite (.txt), Proxyman (.xml), HAR (.har)\\
\textbf{Responsabilité :} Lit le fichier de trafic et détecte des patterns dangereux via regex.
\end{tcolorbox}

\begin{center}
\renewcommand{\arraystretch}{1.3}
\begin{tabular}{lp{8.5cm}}
\toprule
\textbf{Méthode} & \textbf{Ce qu'elle détecte} \\
\midrule
\rowcolor{primary!5}
\texttt{find\_bearer\_tokens()} & Tokens JWT dans les headers \texttt{Authorization: Bearer} \\
\texttt{check\_insecure\_cookies()} & Cookies \texttt{Set-Cookie} sans les flags \texttt{Secure} et \texttt{HttpOnly} \\
\rowcolor{primary!5}
\texttt{find\_secrets\_in\_url()} & Paramètres GET contenant \texttt{token}, \texttt{session}, \texttt{api\_key} \\
\bottomrule
\end{tabular}
\end{center}

\subsubsection{\file{jwt\_parser.py} — Audit des Tokens JWT}

\begin{tcolorbox}[toolbox, title={Rôle : Décodage et analyse de sécurité des JWTs}]
\textbf{Technologie :} PyJWT 2.8+\\
\textbf{Responsabilité :} Décode sans vérifier la signature (mode audit) et analyse les claims.
\end{tcolorbox}

Analyses effectuées :

\begin{itemize}[leftmargin=1.5cm]
  \item \textbf{Algorithme \texttt{none}} : signature totalement désactivée
    $\rightarrow$ \mavscs{MASVS-AUTH-2} — \critical{CRITIQUE}
  \item \textbf{Claim \texttt{exp} absent} : token valide indéfiniment
    $\rightarrow$ \mavscs{MASVS-AUTH-5} — \critical{CRITIQUE}
  \item \textbf{Claim \texttt{sub} absent} : identification du sujet manquante
  \item \textbf{PII dans les claims} : détection de champs \texttt{email}, \texttt{phone}, \texttt{ssn}
\end{itemize}

%% ----------------------------------------------------------
\subsection{Analyseurs — \file{backend/analyzers/}}
%% ----------------------------------------------------------

\subsubsection{\file{semgrep\_runner.py} — Analyse Statique SAST}

\begin{tcolorbox}[toolbox, title={Rôle : Détection de patterns de sécurité dans le code Java/Kotlin décompilé}]
\textbf{Technologie :} Semgrep CLI (outil open source de SAST)\\
\textbf{Responsabilité :} Exécute les règles YAML custom sur le code Java produit par JADX.
\end{tcolorbox}

\textbf{Règles SAST implémentées} (\file{rules/semgrep/auth.yaml}) :

\begin{center}
\renewcommand{\arraystretch}{1.3}
\begin{tabular}{lll}
\toprule
\textbf{Règle} & \textbf{Ce qu'elle détecte} & \textbf{MASVS} \\
\midrule
\rowcolor{red!6}
\texttt{hardcoded-api-key} & Secrets codés en dur (token, API key, JWT) & \mavscs{MASVS-AUTH-7} \\
\texttt{jwt-algorithm-none} & Algorithme JWT \texttt{none} (pas de signature) & \mavscs{MASVS-AUTH-2} \\
\rowcolor{red!6}
\texttt{insecure-shared-preferences} & \texttt{getSharedPreferences()} sans chiffrement & \mavscs{MASVS-AUTH-7} \\
\texttt{token-in-url-parameter} & Token en paramètre GET (\texttt{?token=}) & \mavscs{MASVS-AUTH-9} \\
\rowcolor{red!6}
\texttt{ssl-verification-disabled} & TrustAllCertificates / verification MITM & \mavscs{MASVS-NETWORK-1} \\
\texttt{sensitive-data-in-logs} & \texttt{Log.d()} avec token/password & \mavscs{MASVS-AUTH-7} \\
\rowcolor{red!6}
\texttt{cleartext-http-usage} & URLs HTTP en clair (\texttt{http://}) & \mavscs{MASVS-NETWORK-1} \\
\bottomrule
\end{tabular}
\end{center}

\subsubsection{\file{flow\_analyzer.py} — Analyste du Trafic (DAST)}

\begin{tcolorbox}[toolbox, title={Rôle : Détection de problèmes de sécurité dans le trafic HTTP capturé}]
\textbf{Dépendance :} Utilise TrafficParser comme source de données\\
\textbf{Responsabilité :} Transforme les résultats bruts du TrafficParser en objets Finding MASVS.
\end{tcolorbox}

Détections réalisées par le FlowAnalyzer :

\begin{itemize}[leftmargin=1.5cm]
  \item \textbf{Fuite de token HTTP} : Bearer token présent dans les logs
    $\rightarrow$ \mavscs{MASVS-AUTH-1}
  \item \textbf{Cookie non sécurisé} : manque \texttt{Secure} ou \texttt{HttpOnly}
    $\rightarrow$ \mavscs{MASVS-AUTH-3}
  \item \textbf{Secret dans l'URL} : paramètres GET sensibles
    $\rightarrow$ \mavscs{MASVS-AUTH-2}
\end{itemize}

\subsubsection{\file{runner.py} — Chef d'Orchestre du Pipeline d'Analyse}

\begin{tcolorbox}[toolbox, title={Rôle : Coordination complète du pipeline SAST/DAST/IA}]
\textbf{Pattern :} Async Background Task (FastAPI BackgroundTasks)\\
\textbf{Responsabilité :} Orchestre séquentiellement toutes les étapes d'analyse et
met à jour le job en temps réel.
\end{tcolorbox}

\textbf{Voir Section \ref{sec:pipeline} pour le détail complet du pipeline d'analyse.}

%% ----------------------------------------------------------
\subsection{Intelligence Artificielle — \file{backend/ai/}}
%% ----------------------------------------------------------

\subsubsection{\file{rag\_engine.py} — Moteur RAG (Retrieval-Augmented Generation)}

\begin{tcolorbox}[toolbox, title={Rôle : Base de Connaissances Vectorielle MASVS}]
\textbf{Technologies :} ChromaDB (base vectorielle), SentenceTransformer (modèle MiniLM)\\
\textbf{Responsabilité :} Indexe les 17 règles MASVS et récupère les plus pertinentes
pour contextualiser les prompts envoyés à Gemini.
\end{tcolorbox}

\textbf{Fonctionnement du RAG en 3 étapes :}

\begin{enumerate}[leftmargin=1.8cm, itemsep=4pt]
  \item \textbf{Indexation} (\textit{une seule fois}) : les 17 règles MASVS du fichier
    \file{masvs\_auth\_full.json} sont transformées en vecteurs numériques (embeddings)
    par le modèle \texttt{all-MiniLM-L6-v2} et stockées dans ChromaDB.

  \item \textbf{Requête} : à chaque génération de checklist, une requête en langage naturel
    est vectorisée et comparée aux vecteurs stockés par similarité cosinus.

  \item \textbf{Contexte} : les $n$ règles les plus similaires (défaut : 5) sont retournées
    comme contexte textuel et injectées dans le prompt Gemini.
\end{enumerate}

\textbf{Avantage :} L'IA ne connaît pas seulement les règles MASVS de manière générale
-- elle reçoit les règles \textit{précisément pertinentes} pour ce type d'authentification
(JWT vs OAuth2 vs Sessions), ce qui améliore considérablement la qualité des checklists.

\subsubsection{\file{llm\_client.py} — Client IA Gemini}

\begin{tcolorbox}[toolbox, title={Rôle : Interface unifiée avec l'IA Google Gemini}]
\textbf{Technologies :} LangChain, langchain-google-genai, Gemini 2.5 Flash\\
\textbf{Responsabilité :} Deux fonctions principales de génération par IA.
\end{tcolorbox}

\begin{description}[leftmargin=2cm, style=nextline]
  \item[\texttt{generate\_executive\_summary()}]
    Génère un résumé exécutif en français des vulnérabilités détectées par le SAST/DAST,
    enrichi par le contexte MASVS du RAG.

  \item[\texttt{generate\_auth\_checklist()}]
    Génère un document JSON structuré contenant :
    \begin{itemize}
      \item Une \textbf{checklist} de tests (SAST + DAST) avec MASVS ref, MASTG test, niveau de risque
      \item Des \textbf{Security Acceptance Criteria} (SAC) en style Gherkin
      \item Des \textbf{scénarios Gherkin} (Given/When/Then) pour User Stories Jira
      \item Des \textbf{risques de conception} (Session Fixation, Token Rotation, etc.)
    \end{itemize}
\end{description}

\textbf{Prompt Engineering :} Le prompt injecte le contexte architectural (type auth,
stockage, timeout, MFA, cert pinning) et le contexte MASVS RAG pour des résultats
hautement contextualisés.

%% ----------------------------------------------------------
\subsection{API REST — \file{backend/api/routes/}}
%% ----------------------------------------------------------

\begin{center}
\renewcommand{\arraystretch}{1.4}
\begin{tabular}{llp{6cm}}
\toprule
\textbf{Route} & \textbf{Méthode} & \textbf{Rôle} \\
\midrule
\rowcolor{primary!5}
\texttt{/api/upload/apk} & POST & Upload d'un APK ou fichier de logs. Retourne \texttt{file\_id} + SHA-256 \\
\texttt{/api/analyze/start} & POST & Lance le pipeline d'analyse en arrière-plan. Retourne \texttt{job\_id} \\
\rowcolor{primary!5}
\texttt{/api/jobs/} & GET & Liste les 20 derniers jobs (tri chronologique) \\
\texttt{/api/jobs/\{id\}} & GET & État en temps réel d'un job (\texttt{pending/running/done/failed}) \\
\rowcolor{primary!5}
\texttt{/api/report/} & GET & Liste des 20 derniers rapports \\
\texttt{/api/report/\{id\}} & GET & Rapport complet (Pydantic AuditReport) \\
\rowcolor{primary!5}
\texttt{/api/export/\{id\}/markdown} & GET & Téléchargement rapport format Markdown \\
\texttt{/api/export/\{id\}/json} & GET & Téléchargement rapport format JSON \\
\rowcolor{primary!5}
\texttt{/api/auth-audit/generate-checklist} & POST & Génération checklist MASVS + SAC par IA \\
\texttt{/health} & GET & Vérification de disponibilité de l'API \\
\bottomrule
\end{tabular}
\end{center}

%% ----------------------------------------------------------
\subsection{Frontend — \file{frontend/src/}}
%% ----------------------------------------------------------

\subsubsection{Architecture des Composants React}

\begin{tcolorbox}[toolbox, title={Stack Frontend}]
\textbf{Technologies :} React 18, TypeScript, TailwindCSS, Vite 5, Lucide React\\
\textbf{Pattern :} Single Page Application (SPA) avec navigation par onglets
\end{tcolorbox}

\begin{description}[leftmargin=2cm, style=nextline]
  \item[\file{App.tsx}]
    Composant racine. Gère la navigation entre 4 onglets via un état local simple.
    Affiche le header glassmorphism avec le statut de l'API.

  \item[\file{Scanner.tsx}]
    Interface principale d'analyse. Supporte le drag-and-drop de fichiers.
    Implémente le \textbf{long-polling} (intervalle 1,5s) pour suivre l'état du job
    en temps réel via \texttt{/api/jobs/\{id\}}.
    Affiche les \texttt{FindingCard} expansibles avec évidence, remédiation et MASTG ref.

  \item[\file{AuthAudit.tsx}]
    Formulaire d'audit architectural complet. Expose les 8 champs de \texttt{AuthAuditRequest}
    (type auth, stockage, refresh token, timeout, logout endpoint, cert pinning, MFA, plateforme).
    Affiche les résultats IA en 4 onglets : \textit{Checklist}, \textit{SAC Agile},
    \textit{Gherkin}, \textit{Risques Design}.
    Permet l'export Markdown et la copie JSON.

  \item[\file{MASVSMapping.tsx}]
    Cartographie statique des 35 contrôles MASVS v2 répartis en 6 catégories
    (AUTH, STORAGE, CRYPTO, NETWORK, RESILIENCE, PLATFORM) avec indicateur de couverture
    par pourcentage.

  \item[\file{ReportHistory.tsx}]
    Tableau de bord de l'historique. Affiche en parallèle les jobs récents (avec badge
    de statut) et les rapports disponibles pour consultation et téléchargement.
\end{description}

%% ==========================================================
\section{Pipeline d'Analyse Complet}
\label{sec:pipeline}
%% ==========================================================

\subsection{Vue d'Ensemble du Flux}

\begin{figure}[ht]
\centering
\begin{tikzpicture}[
  node distance=0.5cm,
  step/.style={draw=#1, fill=#1!10, rounded corners=5pt,
               text width=5.5cm, align=center, minimum height=0.85cm,
               font=\small\bfseries, drop shadow={opacity=0.2}},
  arrow/.style={-Stealth, thick, color=mutedtxt},
  status/.style={font=\scriptsize\ttfamily, color=mutedtxt},
  note/.style={font=\scriptsize\itshape, color=mutedtxt, text width=4cm, align=left}
]

\node[step=primary] (s1) {1. Upload du fichier\\{\footnotesize APK / TXT / XML / HAR}};
\node[step=orange, below=0.4cm of s1] (s2) {2. Création du Job DB\\{\footnotesize status: pending}};
\node[step=purple, below=0.4cm of s2] (s3) {3. Lancement Background Task\\{\footnotesize FastAPI BackgroundTasks}};
\node[step=orange, below=0.4cm of s3] (s4) {4. Détection du type de fichier\\{\footnotesize APK → SAST | Logs → DAST}};

%% Branche SAST
\node[step=red, below left=0.6cm and 2.5cm of s4] (s5a) {5a. SAST Pipeline\\{\footnotesize Apktool → JADX}};
\node[step=red, below=0.4cm of s5a] (s6a) {6a. ManifestParser\\{\footnotesize Activities exportées, permissions}};
\node[step=red, below=0.4cm of s6a] (s7a) {7a. Semgrep SAST\\{\footnotesize 7 règles auth.yaml}};

%% Branche DAST
\node[step=green, below right=0.6cm and 2.5cm of s4] (s5b) {5b. DAST Pipeline\\{\footnotesize TrafficParser}};
\node[step=green, below=0.4cm of s5b] (s6b) {6b. FlowAnalyzer\\{\footnotesize Tokens, cookies, secrets}};
\node[step=green, below=0.4cm of s6b] (s7b) {7b. JWTParser\\{\footnotesize alg=none, exp absent, PII}};

%% Fusion
\node[step=purple, below=2.8cm of s4] (s8) {8. Agrégation des Findings\\{\footnotesize Liste unifiée Finding[]}};
\node[step=primary, below=0.4cm of s8] (s9) {9. RAG + Gemini\\{\footnotesize Résumé exécutif IA}};
\node[step=green, below=0.4cm of s9] (s10) {10. Sauvegarde rapport\\{\footnotesize DB SQLite + JSON}};
\node[step=green, below=0.4cm of s10] (s11) {11. Job status: DONE\\{\footnotesize report\_id disponible}};

%% Flèches principales
\draw[arrow] (s1) -- (s2);
\draw[arrow] (s2) -- (s3);
\draw[arrow] (s3) -- (s4);
\draw[arrow] (s4.south west) to[bend right=5] (s5a.north);
\draw[arrow] (s4.south east) to[bend left=5] (s5b.north);
\draw[arrow] (s5a) -- (s6a);
\draw[arrow] (s6a) -- (s7a);
\draw[arrow] (s5b) -- (s6b);
\draw[arrow] (s6b) -- (s7b);
\draw[arrow] (s7a.south) to[bend left=5] (s8.north west);
\draw[arrow] (s7b.south) to[bend right=5] (s8.north east);
\draw[arrow] (s8) -- (s9);
\draw[arrow] (s9) -- (s10);
\draw[arrow] (s10) -- (s11);

\end{tikzpicture}
\caption{Pipeline d'analyse complet — de l'upload au rapport final}
\end{figure}

\subsection{Détail Étape par Étape}

\begin{description}[leftmargin=1.2cm, style=nextline, itemsep=8pt]

  \item[\textbf{Étape 1 — Upload} (\file{upload.py})]
    L'utilisateur soumet un fichier via l'interface drag-and-drop.
    Le fichier est sauvegardé avec un UUID unique comme préfixe.
    Un hash SHA-256 est calculé en streaming pour identifier l'APK de manière unique.
    \textit{Sortie :} \texttt{\{file\_id, filename, hash, status\}}

  \item[\textbf{Étape 2 — Création du Job} (\file{analyze.py})]
    Un enregistrement \texttt{AnalysisJob} est créé en base de données avec le statut
    \texttt{pending}. Cela permet au frontend de suivre l'état avant même que
    l'analyse démarre.

  \item[\textbf{Étape 3 — Lancement Asynchrone} (\file{analyze.py})]
    La fonction \texttt{run\_full\_analysis()} est enregistrée comme tâche de fond
    (\texttt{BackgroundTasks.add\_task}). La réponse HTTP est immédiatement retournée
    au client avec le \texttt{job\_id} sans attendre la fin de l'analyse.

  \item[\textbf{Étape 4 — Détection du type} (\file{runner.py})]
    La détection se fait par l'extension du fichier :
    \texttt{.apk} → Pipeline SAST ; \texttt{.txt/.xml/.har} → Pipeline DAST.

  \item[\textbf{Étapes 5a→7a — Pipeline SAST} (APK)]
    \begin{enumerate}
      \item \tool{Apktool} décode les ressources XML et le bytecode smali
      \item \tool{ManifestParser} analyse \texttt{AndroidManifest.xml}
      \item \tool{JADX} décompile le bytecode en Java
      \item \tool{Semgrep} applique les 7 règles YAML sur le code Java
    \end{enumerate}

  \item[\textbf{Étapes 5b→7b — Pipeline DAST} (Logs)]
    \begin{enumerate}
      \item \tool{TrafficParser} charge et parse le fichier de logs
      \item \tool{FlowAnalyzer} détecte les patterns HTTP dangereux
      \item \tool{JWTParser} décode et audite chaque token JWT trouvé
    \end{enumerate}

  \item[\textbf{Étape 8 — Agrégation}]
    Tous les \texttt{Finding}s SAST et DAST sont réunis dans une liste unique.
    Le score de sévérité de cada finding est automatiquement calculé.

  \item[\textbf{Étape 9 — IA Gemini + RAG}]
    Le \texttt{RAGEngine} récupère les règles MASVS pertinentes.
    Le \texttt{LLMClient} envoie un prompt enrichi à Gemini Flash pour générer
    un résumé exécutif en français adapté aux findings détectés.

  \item[\textbf{Étapes 10→11 — Persistance et Notification}]
    Le rapport complet est sérialisé en JSON et sauvegardé en base.
    Le job passe au statut \texttt{done} avec le \texttt{report\_id}.
    Le frontend (en polling) détecte le changement et affiche le rapport.

\end{description}

%% ==========================================================
\section{Fonctionnement de l'Audit IA (Auth Audit)}
%% ==========================================================

\subsection{Flux de Génération de Checklist}

L'onglet \textbf{Audit Auth IA} implémente un flux distinct du pipeline d'analyse.
Il est basé sur une interaction \textit{synchrone} avec Gemini :

\begin{figure}[ht]
\centering
\begin{tikzpicture}[
  node distance=0.7cm,
  box/.style={draw=#1!60, fill=#1!8, rounded corners=5pt,
              text width=4.5cm, align=center, font=\small, minimum height=0.8cm},
  arrow/.style={-Stealth, thick, color=primary!70}
]
\node[box=primary] (user) {\textbf{Utilisateur}\\\footnotesize Sélectionne arch. (JWT/OAuth2/Sessions + options)};
\node[box=purple, right=1.5cm of user] (api) {\textbf{API}\\\texttt{/auth-audit/generate-checklist}};
\node[box=orange, right=1.5cm of api] (rag) {\textbf{RAG Engine}\\\footnotesize Recherche MASVS similaires};
\node[box=green, below=0.8cm of rag] (gemini) {\textbf{Gemini Flash}\\\footnotesize Génère JSON structuré};
\node[box=primary, below=0.8cm of user] (result) {\textbf{Frontend}\\\footnotesize 4 onglets : Checklist, SAC, Gherkin, Risques};

\draw[arrow] (user) -- (api);
\draw[arrow] (api) -- (rag);
\draw[arrow] (rag) -- (gemini);
\draw[arrow] (gemini.west) to[bend right=20] (api.south);
\draw[arrow] (api.south) to[bend right=10] (result.north east);
\end{tikzpicture}
\caption{Flux de génération de checklist par l'IA}
\end{figure}

\subsection{Structure de la Requête d'Audit}

Le frontend envoie un objet \texttt{AuthAuditRequest} contenant :

\begin{lstlisting}[language=Python, caption={Paramètres de l'AuthAuditRequest (auth\_audit.py)}]
class AuthAuditRequest(BaseModel):
    auth_type: str                  # JWT | OAuth2 | Sessions
    token_storage: Optional[str]    # EncryptedSharedPrefs | Keystore | Cookie | Memory
    has_refresh_token: Optional[bool]
    logout_endpoint: Optional[str]  # ex: /api/auth/logout
    session_timeout_minutes: Optional[int]
    use_certificate_pinning: Optional[bool]
    use_mfa: Optional[bool]
    platforms: Optional[str] = "Android"  # Android | iOS | Both
\end{lstlisting}

\subsection{Structure de la Réponse JSON}

Gemini retourne un JSON structuré en 4 sections :

\begin{lstlisting}[language=json, caption={Structure du JSON retourné par Gemini}]
{
  "checklist": [
    {
      "id": 1,
      "title": "Verifier la duree de vie du token",
      "description": "Description technique detaillee...",
      "test_type": "BOTH",      // SAST | DAST | BOTH
      "masvs_ref": "MASVS-AUTH-5",
      "mastg_test": "MASTG-TEST-0024",
      "risk_level": "HIGH"      // HIGH | MEDIUM | LOW
    }
  ],
  "acceptance_criteria": [
    "Given un utilisateur authentifie, When le token expire, Then..."
  ],
  "gherkin_scenarios": [
    {
      "user_story": "US-SEC-01 : Gestion du cycle de vie des tokens",
      "scenario": "Given ... \n  When ... \n  Then ...",
      "masvs_ref": "MASVS-AUTH-5"
    }
  ],
  "design_risks": [
    {
      "risk": "Session Fixation",
      "description": "Description du risque de conception",
      "mitigation": "Recommandation concrete"
    }
  ]
}
\end{lstlisting}

%% ==========================================================
\section{Standards de Sécurité — OWASP MASVS v2}
%% ==========================================================

\subsection{Couverture des Contrôles MASVS}

L'outil couvre 35 contrôles répartis en 6 catégories MASVS v2 :

\begin{center}
\begin{tikzpicture}
\begin{axis}[
  xbar,
  width=12cm,
  height=7cm,
  xmin=0, xmax=115,
  xlabel={Couverture (\%)},
  symbolic y coords={
    MASVS-PLATFORM, MASVS-RESILIENCE, MASVS-NETWORK,
    MASVS-CRYPTO, MASVS-STORAGE, MASVS-AUTH
  },
  ytick=data,
  ytick style={draw=none},
  yticklabel style={font=\small\bfseries},
  nodes near coords,
  nodes near coords style={font=\scriptsize, /pgf/number format/fixed},
  bar width=0.45cm,
  enlarge y limits=0.15,
  grid=major,
  grid style={dashed, color=mutedtxt!30},
  axis line style={color=mutedtxt!50},
  tick style={color=mutedtxt!50},
  legend style={at={(0.98,0.02)}, anchor=south east, font=\small}
]
\addplot[fill=primary!70, draw=primary] coordinates {
  (67,MASVS-PLATFORM)
  (33,MASVS-RESILIENCE)
  (100,MASVS-NETWORK)
  (80,MASVS-CRYPTO)
  (67,MASVS-STORAGE)
  (82,MASVS-AUTH)
};
\end{axis}
\end{tikzpicture}
\end{center}

\subsection{Base de Connaissances MASVS (RAG)}

La base de connaissances vectorielle comprend 17 règles :

\begin{longtable}{lp{5cm}p{5.5cm}}
\toprule
\textbf{ID MASVS} & \textbf{Titre} & \textbf{Remédiation clé} \\
\midrule
\endhead
\rowcolor{primary!5}
\mavscs{MASVS-AUTH-1} & Authentification serveur robuste & Tout contrôle côté API, pas mobile \\
\mavscs{MASVS-AUTH-2} & Prévention bypass d'auth & Vérifications serveur uniquement \\
\rowcolor{primary!5}
\mavscs{MASVS-AUTH-3} & Non-usage identifiants falsifiables & UUID v4 + Keystore \\
\mavscs{MASVS-AUTH-4} & Rate Limiting / Bruteforce & 5 tentatives / 15 min max \\
\rowcolor{primary!5}
\mavscs{MASVS-AUTH-5} & Timeout et inactivité & Access: 15 min, Refresh: 30 jours \\
\mavscs{MASVS-AUTH-6} & Invalidation session au logout & Deny List côté serveur pour JWT \\
\rowcolor{primary!5}
\mavscs{MASVS-AUTH-7} & Stockage sécurisé des secrets & EncryptedSharedPrefs + Keystore \\
\mavscs{MASVS-AUTH-8} & Biométrie robuste & BiometricPrompt + opération crypto \\
\rowcolor{primary!5}
\mavscs{MASVS-AUTH-9} & Transfert réseau sécurisé & Authorization: Bearer, pas GET \\
\mavscs{MASVS-AUTH-10} & Rotation des Refresh Tokens & Invalider l'ancien à chaque usage \\
\rowcolor{primary!5}
\mavscs{MASVS-AUTH-11} & Protection fixation de session & Régénérer ID après élévation \\
\mavscs{MASVS-STORAGE-1} & Données sensibles non en clair & AES-256-GCM + EncryptedSharedPrefs \\
\rowcolor{primary!5}
\mavscs{MASVS-STORAGE-2} & Pas de données dans logs & Supprimer Log.d() en production \\
\mavscs{MASVS-STORAGE-3} & Pas de données dans backups & android:allowBackup=false \\
\rowcolor{primary!5}
\mavscs{MASVS-NETWORK-1} & TLS activé et configuré & TLS 1.3, rejeter les certificats invalides \\
\mavscs{MASVS-NETWORK-2} & Certificate Pinning & SPKI pinning via OkHttp \\
\rowcolor{primary!5}
\mavscs{MASVS-RESILIENCE-1} & Anti-tampering & R8/ProGuard, RootBeer, intégrité APK \\
\bottomrule
\caption{Base de connaissances MASVS indexée dans ChromaDB}
\end{longtable}

\newcommand{\mavscs}[1]{\textcolor{mavscolor}{\texttt{\small#1}}}

%% ==========================================================
\section{Flux de Travail DevSecOps (Shift-Left)}
%% ==========================================================

\subsection{Intégration dans un Processus Agile}

L'outil produit des livrables directement utilisables dans un workflow Agile/Scrum :

\begin{figure}[ht]
\centering
\begin{tikzpicture}[
  node distance=0.8cm,
  phase/.style={draw=#1!60, fill=#1!10, rounded corners=6pt,
                text width=3.2cm, align=center, minimum height=1.2cm,
                font=\small\bfseries},
  deliverable/.style={draw=mutedtxt!60, fill=white, rounded corners=3pt,
                      text width=2.8cm, align=center, font=\scriptsize,
                      minimum height=0.7cm, dashed},
  arrow/.style={-Stealth, thick, color=mutedtxt}
]

\node[phase=primary] (conception) {\textbf{Conception}\\ \footnotesize Sprint Planning};
\node[phase=orange, right=1cm of conception] (dev) {\textbf{Développement}\\ \footnotesize Coding / PR};
\node[phase=purple, right=1cm of dev] (test) {\textbf{Test}\\ \footnotesize QA / AppSec};
\node[phase=green, right=1cm of test] (release) {\textbf{Release}\\ \footnotesize Déploiement};

\node[deliverable, below=0.7cm of conception] (d1) {Checklists MASVS\\+ Risques Design};
\node[deliverable, below=0.7cm of dev] (d2) {SAC Agile\\(User Stories)};
\node[deliverable, below=0.7cm of test] (d3) {Scénarios Gherkin\\SAST + DAST};
\node[deliverable, below=0.7cm of release] (d4) {Rapport PDF/MD\\+ JSON export};

\draw[arrow] (conception) -- (dev);
\draw[arrow] (dev) -- (test);
\draw[arrow] (test) -- (release);

\draw[->, dashed, color=primary!60] (conception) -- (d1);
\draw[->, dashed, color=orange!60] (dev) -- (d2);
\draw[->, dashed, color=purple!60] (test) -- (d3);
\draw[->, dashed, color=green!60] (release) -- (d4);

\end{tikzpicture}
\caption{Intégration des livrables de l'outil dans le cycle Agile}
\end{figure}

\subsection{Exemple de Security Acceptance Criteria (SAC)}

\begin{tcolorbox}[mavsbox, title={User Story Jira — US-SEC-05 : Invalidation de session}]
\textbf{En tant que} système d'authentification,\\
\textbf{Je veux} invalider la session côté serveur lors d'une déconnexion,\\
\textbf{Afin que} les tokens volés après logout ne puissent plus être utilisés.\\[6pt]
\textbf{Critères d'acceptation :}
\begin{itemize}
  \item \textit{GIVEN} un utilisateur avec un token JWT valide
  \item \textit{WHEN} il appelle \texttt{POST /api/auth/logout}
  \item \textit{THEN} le token est ajouté à la Deny List serveur
  \item \textit{AND} toute requête ultérieure avec ce token retourne HTTP 401
  \item \textit{AND} le log d'audit enregistre l'événement de déconnexion
\end{itemize}
\textbf{Référence MASVS :} \mavscs{MASVS-AUTH-6} | Test MASTG : \texttt{MASTG-TEST-0025}
\end{tcolorbox}

%% ==========================================================
\section{Guide de Déploiement}
%% ==========================================================

\subsection{Prérequis Système}

\begin{tcolorbox}[warnbox, title={Outils Externes Requis}]
\begin{itemize}[itemsep=2pt]
  \item \textbf{Python 3.11+} avec environnement virtuel (\texttt{.venv})
  \item \textbf{Node.js 18+} pour le frontend (npm)
  \item \textbf{Apktool} : outil de décompilation APK (Java requis)
  \item \textbf{JADX} : décompilateur APK vers Java
  \item \textbf{Semgrep} : outil SAST (Linux/Mac natif ; WSL recommandé sur Windows)
  \item \textbf{Docker} (optionnel) : pour le déploiement conteneurisé
  \item \textbf{MobSF} (optionnel) : sur port \textbf{8008} (pas 8000)
\end{itemize}
\end{tcolorbox}

\subsection{Configuration de l'Environnement}

\begin{lstlisting}[language=bash, caption={Installation et démarrage du projet}]
# 1. Cloner et configurer
cp .env.example .env
# Editer .env : ajouter GEMINI_API_KEY

# 2. Backend
python -m venv .venv
.venv\Scripts\activate          # Windows
pip install -r backend/requirements.txt

# Demarrer le backend (port 8000)
python -m uvicorn backend.main:app --reload --host 0.0.0.0 --port 8000

# 3. Frontend (dans un autre terminal)
cd frontend
npm install
npm run dev                     # http://localhost:5173

# 4. MobSF (optionnel — port 8008)
docker run -it -p 8008:8000 opensecurity/mobile-security-framework-mobsf:latest
# Recuperer l'API Key sur http://localhost:8008
# Puis : MOBSF_URL=http://localhost:8008 dans .env
\end{lstlisting}

\subsection{Variables d'Environnement Complètes}

\begin{center}
\renewcommand{\arraystretch}{1.4}
\begin{tabular}{llp{5cm}}
\toprule
\textbf{Variable} & \textbf{Requis} & \textbf{Description} \\
\midrule
\rowcolor{red!6}
\texttt{GEMINI\_API\_KEY} & \textbf{Oui} & Clé Google AI Studio (gratuit sur aistudio.google.com) \\
\texttt{DATABASE\_URL} & Non & SQLite par défaut (\texttt{sqlite:///./analyzer.db}) \\
\rowcolor{primary!5}
\texttt{SECRET\_KEY} & Prod & Clé JWT backend (32 hex chars via openssl) \\
\texttt{MOBSF\_URL} & Non & \texttt{http://localhost:8008} (Docker MobSF) \\
\rowcolor{primary!5}
\texttt{MOBSF\_API\_KEY} & Non & Clé depuis dashboard MobSF (\texttt{localhost:8008}) \\
\texttt{UPLOAD\_DIR} & Non & Auto-calculé relatif au projet \\
\rowcolor{primary!5}
\texttt{CHROMA\_DIR} & Non & Auto-calculé (\texttt{data/chroma\_db/}) \\
\bottomrule
\end{tabular}
\end{center}

%% ==========================================================
\section{Sécurité et Bonnes Pratiques Implémentées}
%% ==========================================================

\subsection{Sécurité du Backend}

\begin{tcolorbox}[successbox, title={Mesures de Sécurité Implémentées}]
\begin{itemize}[itemsep=4pt]
  \item \textbf{Validation des entrées :} Pydantic v2 valide automatiquement tous les
    body de requêtes API avec des contraintes strictes (types, formats, ranges)
  \item \textbf{Isolation des sessions DB :} Chaque opération base de données utilise
    sa propre session \texttt{SessionLocal()} fermed en \texttt{finally}
  \item \textbf{Sanitisation des fichiers :} Les noms de fichiers sont préfixés par UUID
    pour éviter les attaques de path traversal
  \item \textbf{CORS configuré :} En production, restreindre \texttt{allow\_origins}
    à l'URL du frontend uniquement
  \item \textbf{Hashing SHA-256 :} Chaque APK est identifié par son hash cryptographique
    pour détecter les doublons
  \item \textbf{Chemins absolus :} Tous les chemins sont calculés via \texttt{\_\_file\_\_}
    pour éviter les dépendances au répertoire courant
\end{itemize}
\end{tcolorbox}

\subsection{Scoring des Vulnérabilités}

Le score de chaque finding est calculé par la formule :
$$\text{Score} = \text{Impact} \times \text{Exploitability} \times \text{Exposure} \quad \in [1, 27]$$

\begin{center}
\renewcommand{\arraystretch}{1.3}
\begin{tabular}{ccc}
\toprule
\textbf{Score} & \textbf{Sévérité} & \textbf{Pénalité /100} \\
\midrule
\rowcolor{red!10}
$\geq 18$ & \critical{CRITIQUE} & $-15$ points \\
\rowcolor{orange!10}
$\geq 9$ & \textcolor{orange}{\textbf{MAJEUR}} & $-5$ points \\
$< 9$ & \textcolor{green!70!black}{\textbf{MINEUR}} & $-1$ point \\
\bottomrule
\end{tabular}
\end{center}

$$\text{Score Global} = \max\!\left(0,\ 100 - \sum_{\text{findings}} \text{pénalité}(f)\right)$$

%% ==========================================================
\section{Conclusion}
%% ==========================================================

\subsection{Récapitulatif des Capacités}

\begin{center}
\renewcommand{\arraystretch}{1.4}
\begin{tabular}{llc}
\toprule
\textbf{Capacité} & \textbf{Technologie} & \textbf{Statut} \\
\midrule
\rowcolor{green!6}
Analyse SAST Android (APK) & Apktool + JADX + Semgrep & \ok{Opérationnel} \\
Analyse DAST trafic HTTP & TrafficParser + FlowAnalyzer & \ok{Opérationnel} \\
\rowcolor{green!6}
Audit tokens JWT & PyJWT + JWTParser & \ok{Opérationnel} \\
RAG MASVS contextuel & ChromaDB + MiniLM & \ok{Opérationnel} \\
\rowcolor{green!6}
Génération IA checklist & Gemini 2.5 Flash & \ok{Opérationnel} \\
Scénarios Gherkin & Gemini + Prompt Eng. & \ok{Opérationnel} \\
\rowcolor{green!6}
Export Markdown / JSON & FastAPI + Pydantic & \ok{Opérationnel} \\
Suivi temps réel (polling) & FastAPI Jobs API & \ok{Opérationnel} \\
\rowcolor{green!6}
Couverture MASVS v2 & 17 règles RAG + 7 SAST & \ok{Opérationnel} \\
Intégration MobSF & Port 8008 via .env & \textcolor{orange}{\textbf{Optionnel}} \\
\bottomrule
\end{tabular}
\end{center}

\subsection{Architecture : Forces Clés}

\begin{enumerate}[leftmargin=1.8cm, itemsep=6pt]
  \item \textbf{Shift-Left Security :} La sécurité est intégrée dès la phase de conception
    grâce aux checklists et SAC générées avant même d'écrire du code.

  \item \textbf{Pipeline asynchrone :} L'utilisation de \texttt{BackgroundTasks} FastAPI
    permet des analyses longues (décompilation APK, Semgrep) sans bloquer le client.

  \item \textbf{RAG contextualisé :} L'injection de la base de connaissances MASVS via
    ChromaDB garantit que Gemini produit des recommandations précises et standards-alignées.

  \item \textbf{Architecture stratifiée :} Parsers → Analyzers → AI → API → Frontend
    avec séparation nette des responsabilités, facilitant les tests et la maintenance.

  \item \textbf{0 dépendance circulaire :} Chaque module a des dépendances unidirectionnelles
    clairement définies.
\end{enumerate}

\vspace{1cm}

\begin{tcolorbox}[
  enhanced,
  colback=darkbg,
  colframe=primary,
  fontupper=\color{white}\small,
  arc=8pt, boxrule=0.8pt,
  left=10pt, right=10pt, top=8pt, bottom=8pt
]
\textbf{Auth \& Session Security Analyzer} · OWASP MASVS v2 · MASTG · DevSecOps\\
\textcolor{mutedtxt}{Pipeline SAST/DAST + RAG + Gemini API · Python FastAPI + React TypeScript}\\
\textcolor{mutedtxt}{Build : \texttt{\textbf{0 erreurs}} Frontend | 17/17 imports Backend · Avril 2026}
\end{tcolorbox}

\end{document}
